Extended reality is a comprehensive term that refers to all different ways of combining virtual and real entities in a human-machine interface system. It is usually decomposed into four classes (augmented reality, augmented virtuality, mixed reality, and virtual reality) that differ from one another based on the level of reality and virtuality involved in the interface system. 

\citeonline{milgram1994taxonomy} ordered these classes and created the concept of “virtuality continuum”, as illustrated in Figure \ref{fig:virtuality_continuum}. On the left, the real environment represent the cases where the stimuli is not produced by a computer or any other digital system. Along the path to the right, the environment starts to incorporate digital elements until it reaches the far right, where all the elements in the environment are virtual and have a digital origin \cite{nijholt2005virtuality, doolani2020review}.

\begin{figure}[!h]
    \tikzstyle{arrow} = [ccmDBlue, rounded corners, line width = 2mm, ->]
    \tikzstyle{--blue} = [ccmDBlue, rounded corners, line width = 2mm]
    \tikzstyle{--black} = [rounded corners, line width = 1mm]
    
    %\tikzstyle{arrow_flow} = [ccmblue, rounded corners, line width = 2mm, ->]
    %\tikzstyle{arrow_return} = [ccmred, rounded corners, line width = 2mm, ->]
    
    \resizebox{0.80\width}{!}{
    \begin{tikzpicture}[node distance=1cm]
        \centering
    
        \node (left) {};
        
        \node (reality) [right of = left, xshift = 2cm]{\includegraphics[width=.15\textwidth]{Fundamentação/Realidade Extendida/real.png}}
        node(t_reality)[below of = reality,yshift=-0.75cm] {Real}
        node(t_reality2)[below of = t_reality,yshift=0.5cm] {Environment};
        
        \node (midLeft) [right of = reality, xshift = 1cm] {};
        
        \node (ar) [right of=reality, xshift=3cm] {\includegraphics[width=.15\textwidth]{Fundamentação/Realidade Extendida/ar.png}}
        node(t_ar)[below of = ar,yshift=-0.75cm] {Augmented}
        node(t_ar2)[below of = t_ar,yshift=0.5cm] {Reality};
        
        \node (av) [right of=ar, xshift=3cm] {\includegraphics[width=.15\textwidth]{Fundamentação/Realidade Extendida/av.png}}
        node(t_av)[below of = av,yshift=-0.75cm] {Augmented}
        node(t_av2)[below of = t_av,yshift=0.5cm] {Virtuality};
        
        \node (vr) [right of=av, xshift=3cm] {\includegraphics[width=.15\textwidth]{Fundamentação/Realidade Extendida/vr.png}}
        node(t_vr)[below of = vr,yshift = -0.75cm] {Virtual} 
        node(t_vr2)[below of = t_vr,yshift = 0.5cm] {Reality};
        
        \node (midRight) [left of = vr, xshift = -1cm] {};
        
        \node (right) [right of = vr, xshift = 2cm] {};
        
        \node (mr) [above of=ar, right of=ar, xshift=1cm, yshift=3cm] {\includegraphics[width=.15\textwidth]{Fundamentação/Realidade Extendida/mr.png}} 
        node(t_mr) [below of = mr, yshift = -0.50cm]{Mixed}
        node(t_mr2) [below of = t_mr, yshift = 0.5cm]{Reality};
        
        \draw [arrow] (reality) to (left);
        \draw [--blue] (reality) to (ar);
        \draw [--blue] (ar) to (av);
        \draw [--blue] (av) to (vr);
        \draw [arrow] (vr) to (right);
        \draw [--black] (mr.west) -- +(-2.6,0) to (midLeft);
        \draw [--black] (mr.east) -- +(2.6,0) to (midRight);
         
        
        \node (left) [left of = reality, xshift = -2cm] {};
        \node (right) [right of = vr, xshift = 2cm] {};
    
        
    \end{tikzpicture}
    }
    \centering
    \caption{The Virtuality Continuum concept \cite{milgram1994taxonomy}}
    \label{fig:virtuality_continuum}
\end{figure}

The extreme left means full reality, where the stimuli does not come, or is not produced by any computer or any other digital system. Along the path to the right, the environment starts to have some digital elements until it reaches the far right, where all the elements in the environment have a digital origin  The first step from the Real Environment to Virtual Reality is the Augmented Reality.

In an augmented reality system, the user sees some digital elements, which could be text, images, video, etc, that are laid over the real environment without making the user to lose his sense of presence in the real world. Augmented reality can be used to assist workers in manufacturing and assembly tasks, as well as training \cite{doolani2020review, farrell2018learning, ma2007virtuality}.
    
While augmented reality brings digital elements to the real environment, augmented virtuality creates an environment that could only exist digitally, such as a fantasy world from games or movies. This scenario is the background of some other activity that is being done by the user in the real environment. An example is to train a pilot in a virtual environment but with a real mock-up of the cockpit, which provides physical buttons and inceptors for the pilot to touch and hold \cite{farshid2018go}. Another example is to play sports, such as tennis, golf or baseball, in a complete digital arena but using the real equipment with a tracker.

Mixed Reality stays in between the real and virtual environments. Differently from augmented reality and augmented virtuality, in a mixed reality system the user can manipulate digital elements as if they were inside the real world  \cite{doolani2020review}.One example is when a client from a furniture store uses mixed reality not only to see how the furniture fits inside his room, but he can also move it, change its colour, size and shape before buying or even going to the shop.
    
On the far right of the virtuality continuum, the virtual reality is when the user is the only non-digital element, everything else is digital, immersing the user in a virtual environment, but, of course, inside the physical limits of the real environment \cite{ma2007virtuality}. If the feeling of presence of that environment is well done, the user can momentarily forget about the real environment and act and react accordingly to the virtual environment \cite{farrell2018learning}. 
    
Virtual reality is a powerful tool that allows a user to be transported to a tridimensional environment that could be out of reach or that doesn’t exist but is needed for testing or training reasons \cite{mujber2004virtual}. Inside this virtual environment, the user can walk, look around and feel as if the environment was real \cite{salah2019virtual}.
    
The Figure \ref{fig:ar_av_mr} shows the representations of each of these Extended Reality classes.
    
\begin{figure}[!htb]
    \tikzstyle{arrow} = [ccmblue, rounded corners, line width = 2mm, ->]
    \tikzstyle{--blue} = [ccmblue, rounded corners, line width = 2mm]
    \tikzstyle{--black} = [rounded corners, line width = 1mm]
    
    \resizebox{0.85\width}{!}{
    \begin{tikzpicture}[node distance=1cm]
        \centering
    
        \node (ar) {\includegraphics[width=0.3\textwidth]{Fundamentação/Realidade Extendida/ar.png}}
        node(t_ar)[below of = ar,yshift=-2.0cm] {Augmented}
        node(t_ar2)[below of = t_ar,yshift=0.5cm] {Reality};
        
        \node (av) [right of=ar, xshift=5cm, yshift = 0.1cm] {\includegraphics[width=.34\textwidth]{Fundamentação/Realidade Extendida/av.png}}
        node(t_av)[below of = av,xshift = 1.0cm, yshift=-2.0cm] {Augmented}
        node(t_av2)[below of = t_av,yshift=0.5cm] {Virtuality};
        
        \node (mr) [below of=ar, yshift=-5.5cm] {\includegraphics[width=.31\textwidth]{Fundamentação/Realidade Extendida/mr.png}}
        node(t_mr)[below of = mr,yshift=-2.0cm] {Mixed}
        node(t_mr2)[below of = t_mr,yshift=0.5cm] {Reality};
        
        \node (vr) [right of=mr, xshift=5.0cm, yshift = -0.5cm] {\includegraphics[width=.25\textwidth]{Fundamentação/Realidade Extendida/vr.png}}
        node(t_vr)[below of = vr,yshift = -1.5cm] {Virtual} 
        node(t_vr2)[below of = t_vr,yshift = 0.5cm] {Reality};
        
        \node (n) [right of = ar, above of = ar , xshift = 2cm, yshift = 3.0cm] {};
        \node (l) [right of = av, below of = av, xshift = 3cm, yshift = -3.1cm] {};
        \node (s) [left of = vr, below of = vr, xshift = -2cm, yshift = -3.0cm] {};
        \node (o) [left of = mr, above of = mr, xshift = -3cm, yshift = 1.6cm] {};
        
        
        %\draw [arrow] (reality) to (left);
        %\draw [--blue] (reality) to (ar);
        %\draw [--blue] (ar) to (av);
        %\draw [--blue] (av) to (vr);
        %\draw [arrow] (vr) to (right);
        \draw [--black] (n) to (s);
        \draw [--black] (l) to (o);
         
        
        %\node (left) [left of = reality, xshift = -2cm] {};
        %\node (right) [right of = vr, xshift = 2cm] {};
    
        
    \end{tikzpicture}
    }
    \centering
    \caption{A represantion of the differences between AR, AV, MR and VR. Made by the author}
    \label{fig:ar_av_mr}
\end{figure}